\documentclass[11pt]{article}
\usepackage[a4paper,margin=25mm]{geometry}
\usepackage{amsmath,amssymb,amsthm}
\usepackage{booktabs}
\usepackage[hidelinks]{hyperref}
\usepackage{lmodern}
\usepackage{microtype}

\newtheorem{theorem}{Theorem}
\newtheorem{lemma}[theorem]{Lemma}
\newtheorem{proposition}[theorem]{Proposition}
\newtheorem{corollary}[theorem]{Corollary}
\theoremstyle{remark}
\newtheorem{remark}[theorem]{Remark}

\newcommand{\Fp}{\mathbb{F}_p}
\newcommand{\Z}{\mathbb{Z}}
\newcommand{\Hc}{H(c,p)}
\newcommand{\Gp}{G(p)}
\newcommand{\Lc}{\mathcal{L}}
\newcommand{\Pset}{\mathcal{P}}
\newcommand{\cl}[1]{\langle #1\rangle}

\title{The Hall--Jackson--Sudbery--Wild window is tight:\\
a modular hyperbola in the $2p\times 2p$ grid $G(p)$ contains\\ no more than $3(p-1)$ points in general position}
\author{Alex Komang\thanks{Independent researcher, Bali, Indonesia.}\\[2pt]
\small with three autonomous Claude agents (Anthropic);\\ \small this note was written by the third agent --- see the statement at the end}
\date{August 2026 --- draft v0.1}

\begin{document}
\maketitle

\begin{abstract}
Let $p$ be an odd prime, $c\in\Fp^*$, and let $\Hc=\{(x,y)\in\Z^2: xy\equiv c \pmod p\}$ be the modular hyperbola.
Hall, Jackson, Sudbery and Wild (1975) observed that the window
$\Gp=[-\tfrac{p-1}{2},\tfrac{3p-1}{2}]\times[0,2p-1]$ contains exactly four points of each of the $p-1$ residue classes of
$\Hc$, at the vertices of a $p\times p$ square, and that keeping three of the four (the twelve outer blocks $T_2$ of the
$4\times 4$ block subdivision of $\Gp$) yields $3(p-1)$ points with no three collinear.  We prove that this is optimal:
\emph{every} subset of $\Hc\cap\Gp$ with no three collinear points has at most $3(p-1)$ points.  The proof is a
double count over the lines of slope $\pm1$ that carry three or four points of $\Hc\cap\Gp$.  The argument also
determines all maximum sets: their number is $9^{s}$, where $s\in\{0,1,2\}$ is the number of quadratic residues among
$c$ and $-c$; for $s=0$ the HJSW set $\Hc\cap T_2$ is the unique maximum.  Everything is verified by an independent
program for all primes $p\le 101$.
\end{abstract}

\section{Statement}

Throughout, $p$ is an odd prime, $h=(p-1)/2$, $c\in\Fp^*$, and
\[
 \Hc=\{(x,y)\in\Z^2:\ xy\equiv c \pmod p\},\qquad
 \Gp=\{-h,\dots,3h+1\}\times\{0,\dots,2p-1\},
\]
\[
 \Pset:=\Hc\cap\Gp .
\]
$\Gp$ is a $2p\times2p$ grid (this is the grid $G$ of \cite[(2.1)]{KNS25}, taken from \cite{HJSW75}).  A set is
\emph{lawful} (in general position, no-three-in-line) if it contains no three collinear points.

Two points are \emph{congruent} if their coordinates agree modulo $p$; a \emph{class} is a congruence class,
identified with an element of $\Fp\times\Fp$.  Points of $\Hc$ have both coordinates prime to $p$, so the classes of
$\Hc$ are the $p-1$ pairs $\kappa=(a,b)\in\Fp^*\times\Fp^*$ with $ab=c$.  Since each coordinate range of $\Gp$ is an
interval of $2p$ consecutive integers, it contains exactly two representatives of every residue class; hence
every class $\kappa$ of $\Hc$ has exactly four points in $\Gp$ and $|\Pset|=4(p-1)$.  We write
$(x_\kappa,y_\kappa)$ for the \emph{base copy} of $\kappa$, its unique point in
$\{-h,\dots,h\}\times\{1,\dots,p-1\}$ (necessarily $x_\kappa\ne0$), and
\[
 \kappa(r,s):=(x_\kappa+rp,\;y_\kappa+sp),\qquad r,s\in\{0,1\},
\]
for its four points in $\Gp$.  Following \cite{HJSW75,KNS25} we sort the classes into four \emph{types} according
to the position of the base copy:
\[
 A:\ x_\kappa>0,\ y_\kappa>h;\qquad B:\ x_\kappa<0,\ y_\kappa>h;\qquad C:\ x_\kappa>0,\ y_\kappa\le h;\qquad
 D:\ x_\kappa<0,\ y_\kappa\le h .
\]
(In the notation of \cite[Def.~2.3]{KNS25}: $\kappa(0,0)\in A_{0,0}$, $B_{-1,0}$, $C_{0,0}$, $D_{-1,0}$ respectively; the
HJSW set $T_2$ omits, for each type, the copy $A_{0,0}$, $B_{0,0}$, $C_{0,1}$, $D_{0,1}$, i.e.\ $\kappa(0,0)$ for
$A$, $\kappa(1,0)$ for $B$, $\kappa(0,1)$ for $C$, $\kappa(1,1)$ for $D$; the union of these omitted blocks is the middle
block $M$ of \cite[\S2.3]{KNS25}.)  We write $n_A,n_B,n_C,n_D$ for the numbers of classes of each type, so
$n_A+n_B+n_C+n_D=p-1$.  Finally, for a class $\kappa$ put
\[
 d_\kappa:=x_\kappa-y_\kappa,\qquad e_\kappa:=x_\kappa+y_\kappa .
\]

Three involutions act on the classes of $\Hc$:
\[
 \sigma(a,b)=(-b,-a),\qquad \tau(a,b)=(b,a),\qquad \nu(a,b)=(-a,-b)=\sigma\tau(a,b)=\tau\sigma(a,b).
\]
$\sigma$ has fixed classes iff $-c$ is a quadratic residue (then exactly two: $(a,-a)$ with $a^2=-c$), $\tau$ has fixed
classes iff $c$ is a quadratic residue (two: $(a,a)$ with $a^2=c$), and $\nu$ has none.  Put
\[
 f_\sigma:=\#\{\kappa:\sigma\kappa=\kappa\}\in\{0,2\},\qquad f_\tau:=\#\{\kappa:\tau\kappa=\kappa\}\in\{0,2\},
\]
\[
 s:=\tfrac12(f_\sigma+f_\tau)=[\,-c\text{ is a QR}\,]+[\,c\text{ is a QR}\,]\in\{0,1,2\}.
\]

\begin{theorem}\label{thm:main}
Every lawful subset of $\Pset=\Hc\cap\Gp$ has at most $3(p-1)$ points.  The bound is attained, and the number of
lawful subsets of size $3(p-1)$ is exactly $9^{s}$.  In particular, if neither $c$ nor $-c$ is a quadratic residue
modulo $p$, the HJSW set $\Hc\cap T_2$ is the unique maximum lawful subset of $\Pset$.
\end{theorem}

For $c=1$: there are $81$ maximum sets if $p\equiv1\pmod4$ and $9$ if $p\equiv3\pmod4$; all of them agree with the
HJSW set outside the two or four classes $(\pm a,\mp a)$, $(\pm a,\pm a)$ with $a^2=\mp1$.

\section{Lines through three points of $\Pset$}

\begin{lemma}\label{lem:lagrange}
Let $\ell$ be a line containing at least two points of $\Hc$.  Then the points of $\Hc\cap\ell$ lie in at most two
classes.  More precisely, if $\ell$ contains a point of class $\kappa$, then the classes met by $\ell$ are:
only $\kappa$ if $\ell$ has slope $0$ or $\infty$; among $\{\kappa,\sigma\kappa\}$ if $\ell$ has slope $+1$;
among $\{\kappa,\tau\kappa\}$ if $\ell$ has slope $-1$.
\end{lemma}

\begin{proof}
Let $(u,v)\in\Z^2$ be a primitive direction vector of $\ell$ and $(x_0,y_0)\in\Hc\cap\ell$.  The lattice points of
$\ell$ are $(x_0+tu,y_0+tv)$, $t\in\Z$, and such a point lies on $\Hc$ iff
$uv\,t^2+(x_0v+y_0u)\,t\equiv0\pmod p$ (we used $x_0y_0\equiv c$).  The class of the point depends only on
$t\bmod p$.  If $uv\not\equiv0$, the congruence has at most two roots modulo $p$; if exactly one of $u,v$ is
divisible by $p$ (both cannot be, as $\gcd(u,v)=1$), it reduces to $x_0v\,t\equiv0$ or $y_0u\,t\equiv0$ with a unit
coefficient, so $t\equiv0$: one class.  For $(u,v)=(1,0),(0,1)$ this gives the claim for slopes $0,\infty$.
For $(u,v)=(1,1)$ the roots are $t\equiv0$ and $t\equiv-(x_0+y_0)$, the latter giving the point
$(-y_0,-x_0)$ modulo $p$, of class $\sigma\kappa$.  For $(u,v)=(1,-1)$ the roots are $t\equiv0$ and
$t\equiv y_0-x_0$, giving $(y_0,x_0)$, of class $\tau\kappa$.
\end{proof}

\begin{lemma}\label{lem:copies}
Two distinct points of one class in $\Gp$ differ by $\pm(p,0)$, $\pm(0,p)$, $\pm(p,p)$ or $\pm(p,-p)$; the line
through them has slope $0,\infty,+1$ or $-1$ and contains no third point of that class.  The point $\kappa(r,s)$ lies
on the slope-$(+1)$ line $x-y=d_\kappa+(r-s)p$ and on the slope-$(-1)$ line $x+y=e_\kappa+(r+s)p$.
\end{lemma}

\begin{proof}
The four points of a class are the vertices of an axis-parallel square of side $p$; no three are collinear.
The last sentence is a direct computation.
\end{proof}

\begin{corollary}\label{cor:slopes}
A line containing at least three points of $\Pset$ has slope $+1$ or $-1$.
\end{corollary}

\begin{proof}
By Lemma~\ref{lem:lagrange} the three points lie in at most two classes, so two of them are congruent, and by
Lemma~\ref{lem:copies} the line has slope $0,\infty,\pm1$.  A line of slope $0$ or $\infty$ meets a single class
(Lemma~\ref{lem:lagrange}), which has at most two points on it (Lemma~\ref{lem:copies}).
\end{proof}

\begin{lemma}[position of the partner classes]\label{lem:partners}
Let $\kappa$ be a class of $\Hc$.
\begin{enumerate}
\item[(i)] If $\kappa\in A$ then $\sigma\kappa\in A$ and $d_{\sigma\kappa}=d_\kappa$; if $\kappa\in D$ then
 $\sigma\kappa\in D$ and $d_{\sigma\kappa}=d_\kappa$; if $\kappa\in C$ then $\sigma\kappa\in B$ and
 $d_{\sigma\kappa}=d_\kappa-p$; if $\kappa\in B$ then $\sigma\kappa\in C$ and $d_{\sigma\kappa}=d_\kappa+p$.
\item[(ii)] If $\kappa\in B$ then $\tau\kappa\in B$ and $e_{\tau\kappa}=e_\kappa$; if $\kappa\in C$ then
 $\tau\kappa\in C$ and $e_{\tau\kappa}=e_\kappa$; if $\kappa\in A$ then $\tau\kappa\in D$ and
 $e_{\tau\kappa}=e_\kappa-p$; if $\kappa\in D$ then $\tau\kappa\in A$ and $e_{\tau\kappa}=e_\kappa+p$.
\end{enumerate}
In particular $\sigma$ preserves the types $A$, $D$ and interchanges $B$, $C$; $\tau$ preserves $B$, $C$ and
interchanges $A$, $D$; the fixed classes of $\sigma$ lie in $A\cup D$ and those of $\tau$ in $B\cup C$; and
$n_A=n_D$, $n_B=n_C$.
\end{lemma}

\begin{proof}
Write $\kappa=(a,b)$ with base copy $(x_\kappa,y_\kappa)$, so $x_\kappa\equiv a$, $y_\kappa\equiv b$.  The base copy of
$\sigma\kappa=(-b,-a)$ has first coordinate $-y_\kappa$ if $y_\kappa\le h$ and $p-y_\kappa$ if $y_\kappa>h$, and second
coordinate $p-x_\kappa$ if $x_\kappa>0$ and $-x_\kappa$ if $x_\kappa<0$ (these are the representatives of $-b$ in
$[-h,h]$ and of $-a$ in $[1,p-1]$).  Going through the four types:
$\kappa\in A$: $(p-y_\kappa,\,p-x_\kappa)$, again of type $A$, difference $x_\kappa-y_\kappa$;
$\kappa\in D$: $(-y_\kappa,\,-x_\kappa)$, type $D$, difference $x_\kappa-y_\kappa$;
$\kappa\in C$: $(-y_\kappa,\,p-x_\kappa)$, type $B$, difference $x_\kappa-y_\kappa-p$;
$\kappa\in B$: $(p-y_\kappa,\,-x_\kappa)$, type $C$, difference $x_\kappa-y_\kappa+p$.
Similarly, the base copy of $\tau\kappa=(b,a)$ has first coordinate $y_\kappa$ if $y_\kappa\le h$ and $y_\kappa-p$
otherwise, and second coordinate $x_\kappa$ if $x_\kappa>0$ and $x_\kappa+p$ otherwise:
$\kappa\in B$: $(y_\kappa-p,\,x_\kappa+p)$, type $B$, sum $e_\kappa$;
$\kappa\in C$: $(y_\kappa,\,x_\kappa)$, type $C$, sum $e_\kappa$;
$\kappa\in A$: $(y_\kappa-p,\,x_\kappa)$, type $D$, sum $e_\kappa-p$;
$\kappa\in D$: $(y_\kappa,\,x_\kappa+p)$, type $A$, sum $e_\kappa+p$.
The remaining assertions follow ($\nu=\sigma\tau$ is a fixed-point-free involution interchanging $A\leftrightarrow D$
and $B\leftrightarrow C$).
\end{proof}

Let $\Lc$ be the set of lines containing at least three points of $\Pset$.  For a class $\kappa$ let
\begin{align*}
 D_\kappa&:=\{x-y=d_\kappa\}\quad(\text{the diagonal of }\kappa;\ \text{it contains }\kappa(0,0),\kappa(1,1)),\\
 E_\kappa&:=\{x+y=e_\kappa+p\}\quad(\text{the antidiagonal of }\kappa;\ \text{it contains }\kappa(0,1),\kappa(1,0)).
\end{align*}

\begin{proposition}\label{prop:lines}
$\Lc$ consists of the following lines, and no others; for each we list its points in $\Pset$.
\begin{enumerate}
\item[(a)] For $\kappa\in A\cup D$ with $\sigma\kappa\ne\kappa$: the common diagonal $D_\kappa=D_{\sigma\kappa}$, with the
 four points $\kappa(0,0),\kappa(1,1),\sigma\kappa(0,0),\sigma\kappa(1,1)$.
\item[(b)] For $\kappa\in B\cup C$ with $\tau\kappa\ne\kappa$: the common antidiagonal $E_\kappa=E_{\tau\kappa}$, with the
 four points $\kappa(0,1),\kappa(1,0),\tau\kappa(0,1),\tau\kappa(1,0)$.
\item[(c)] For every $\kappa\in C$: $D_\kappa$ with the three points $\kappa(0,0),\kappa(1,1),\sigma\kappa(1,0)$
 ($\sigma\kappa\in B$); for every $\kappa\in B$: $D_\kappa$ with the three points
 $\kappa(0,0),\kappa(1,1),\sigma\kappa(0,1)$ ($\sigma\kappa\in C$).
\item[(d)] For every $\kappa\in A$: $E_\kappa$ with the three points $\kappa(0,1),\kappa(1,0),\tau\kappa(1,1)$
 ($\tau\kappa\in D$); for every $\kappa\in D$: $E_\kappa$ with the three points $\kappa(0,1),\kappa(1,0),\tau\kappa(0,0)$
 ($\tau\kappa\in A$).
\end{enumerate}
Consequently
\begin{equation}\label{eq:L}
 |\Lc|=\frac{n_A+n_D-f_\sigma}{2}+\frac{n_B+n_C-f_\tau}{2}+(n_A+n_D)+(n_B+n_C)
 =\frac32(p-1)-\frac{f_\sigma+f_\tau}{2}.
\end{equation}
Moreover, the points of $\Pset$ lying on \emph{no} line of $\Lc$ are exactly $\kappa(1,1)$ for $\sigma$-fixed
$\kappa\in A$, $\kappa(0,0)$ for $\sigma$-fixed $\kappa\in D$, $\kappa(0,1)$ for $\tau$-fixed $\kappa\in B$ and
$\kappa(1,0)$ for $\tau$-fixed $\kappa\in C$; their number is
\begin{equation}\label{eq:Z}
 Z:=\#\{q\in\Pset:\ q\text{ lies on no line of }\Lc\}=f_\sigma+f_\tau .
\end{equation}
The points lying on \emph{two} lines of $\Lc$ are exactly the ``middle-block'' copies $\kappa(0,0)$ ($\kappa\in A$),
$\kappa(1,0)$ ($\kappa\in B$), $\kappa(0,1)$ ($\kappa\in C$), $\kappa(1,1)$ ($\kappa\in D$) of the classes $\kappa$ whose
partner ($\sigma\kappa$ for $A,D$; $\tau\kappa$ for $B,C$) is different from $\kappa$; no point lies on three lines of
$\Lc$.
\end{proposition}

\begin{proof}
Let $\ell\in\Lc$.  By Corollary~\ref{cor:slopes} its slope is $\pm1$; say $+1$ (the case $-1$ is symmetric, with
$\tau$, $e_\kappa$, $E_\kappa$ in place of $\sigma$, $d_\kappa$, $D_\kappa$).  By Lemma~\ref{lem:lagrange} the points of
$\ell\cap\Pset$ lie in two classes $\kappa,\sigma\kappa$ (a single class carries at most two points on a line), and
one of the classes, say $\kappa$, has two points on $\ell$; by Lemma~\ref{lem:copies} these are $\kappa(0,0),\kappa(1,1)$
and $\ell=D_\kappa$, and $\sigma\kappa\ne\kappa$.  The points of $\sigma\kappa$ on $D_\kappa$ are the
$\sigma\kappa(r,s)$ with $d_{\sigma\kappa}+(r-s)p=d_\kappa$.  By Lemma~\ref{lem:partners}(i): for $\kappa\in A\cup D$,
$r=s$, giving (a); for $\kappa\in C$, $r-s=1$, giving $\sigma\kappa(1,0)$; for $\kappa\in B$, $r-s=-1$, giving
$\sigma\kappa(0,1)$; this is (c).  Conversely each listed line has $\ge3$ points.  Note that for $\kappa\in B\cup C$
the diagonals $D_\kappa$ and $D_{\sigma\kappa}$ are distinct lines, each in $\Lc$, whereas for $\kappa\in A\cup D$ they
coincide.  Counting: (a) gives $(n_A+n_D-f_\sigma)/2$ lines, (b) $(n_B+n_C-f_\tau)/2$, (c) $n_B+n_C$, (d) $n_A+n_D$,
which is \eqref{eq:L}.

For the incidences of a single point, take $\kappa\in A$.  Its points $\kappa(0,1),\kappa(1,0)$ lie on $E_\kappa\in\Lc$
(d), and on no line of type (a)/(c) (their slope-$(+1)$ lines $x-y=d_\kappa\pm p$ carry no second point of
$\kappa$ and, by (i), no point of $\sigma\kappa$).  $\kappa(1,1)$ lies on $D_\kappa$, which is in $\Lc$ iff
$\sigma\kappa\ne\kappa$, and on $x+y=e_\kappa+2p=e_{\tau\kappa}+3p$, which carries no point of $\tau\kappa$: so
$\kappa(1,1)$ is on one line of $\Lc$ if $\sigma\kappa\neq\kappa$ and on none otherwise.  $\kappa(0,0)$ lies on
$D_\kappa$ and on $x+y=e_\kappa=e_{\tau\kappa}+p$, i.e.\ on $E_{\tau\kappa}\in\Lc$: two lines if $\sigma\kappa\ne\kappa$,
one otherwise.  The types $D$, $B$, $C$ are analogous (for $B$: $\kappa(1,0)\in D_{\sigma\kappa}$ since
$d_\kappa+p=d_{\sigma\kappa}$, while $\kappa(0,1)$ lies on $x-y=d_\kappa-p=d_{\sigma\kappa}-2p$, which carries no point of
$\sigma\kappa$; for $C$ the roles of $(1,0),(0,1)$ are exchanged; for $D$: $\kappa(1,1)\in E_{\tau\kappa}$ since
$e_\kappa+2p=e_{\tau\kappa}+p$, and $\kappa(0,0)$ lies on $x+y=e_\kappa=e_{\tau\kappa}-p$).  This gives \eqref{eq:Z}
and the last two assertions.
\end{proof}

\section{Proof of Theorem~\ref{thm:main}}

Let $S\subseteq\Pset$ be lawful.  For $q\in\Pset$ let $m(q)$ be the number of lines of $\Lc$ through $q$.  Every line
of $\Lc$ carries at most two points of $S$, hence
\begin{equation}\label{eq:count}
\begin{aligned}
 |S| \;&=\; \#\{q\in S: m(q)=0\}+\#\{q\in S: m(q)\ge1\}
 \;\le\; Z+\sum_{q\in S}m(q)-\sum_{q\in S}\bigl(m(q)-1\bigr)^{+}\\
 &=\; Z+\sum_{\ell\in\Lc}|S\cap\ell|-\sum_{q\in S}\bigl(m(q)-1\bigr)^{+}
 \;\le\; Z+2|\Lc| ,
\end{aligned}
\end{equation}
and by \eqref{eq:L}, \eqref{eq:Z}
\[
 Z+2|\Lc| \;=\; (f_\sigma+f_\tau)+3(p-1)-(f_\sigma+f_\tau) \;=\; 3(p-1).
\]
This proves the bound.  (Lawfulness of $S$ was used only through the lines of $\Lc$; conversely, since every line
carrying three points of $\Pset$ belongs to $\Lc$, a set $S\subseteq\Pset$ is lawful iff $|S\cap\ell|\le2$ for all
$\ell\in\Lc$.)

\medskip
\noindent\emph{Equality.}  $|S|=3(p-1)$ holds iff all inequalities in \eqref{eq:count} are equalities, i.e.\ iff
(1) $S$ contains all $Z$ points lying on no line of $\Lc$; (2) $|S\cap\ell|=2$ for every $\ell\in\Lc$; (3) no point of $S$
lies on two lines of $\Lc$.  Group the classes into orbits of the Klein four-group $\{1,\sigma,\tau,\nu\}$.  A generic
orbit $\{\kappa,\sigma\kappa,\tau\kappa,\nu\kappa\}$ has four distinct classes, all with distinct partners; the six
lines of $\Lc$ meeting it (Proposition~\ref{prop:lines}) meet no other orbit, its four middle-block copies have $m=2$
and are excluded by (3), and each of the six lines then retains exactly two points with $m=1$, which are forced into
$S$ by (2): on this orbit $S$ coincides with $\Hc\cap T_2$.  A non-generic orbit is $\{\kappa,\nu\kappa\}$ with
$\kappa$ (and $\nu\kappa$) fixed by $\sigma$ or by $\tau$; it meets exactly two lines of $\Lc$, each with three points
of $m=1$, plus two points with $m=0$: conditions (1)--(3) allow $3\cdot3=9$ choices, all lawful.  There are
$s=(f_\sigma+f_\tau)/2$ non-generic orbits, whence $9^{s}$ maximum sets.  \qed

\begin{remark}
(a) The count \eqref{eq:count} is a linear-programming certificate: the dual solution assigns weight $1$ to every
line of $\Lc$ (constraint $\sum_{q\in\ell}x_q\le2$) and weight $1$ to the upper bound $x_q\le1$ of each of the $Z$
unconstrained points; so even the fractional relaxation of the problem has value $3(p-1)$.

(b) The theorem says that no lawful set can use the middle block $M$ profitably: adding fourth copies to the HJSW
set is impossible without removing at least as many points.  It complements Observations 2.4--2.6 of \cite{KNS25}
(which yield the lower bound $3(p-1)$) and settles, for a single hyperbola in its own window, the question raised
there whether the HJSW construction can be improved: within $\Hc\cap\Gp$ it cannot; any improvement must use
points outside $\Hc\cap\Gp$ (another hyperbola, another window, or non-hyperbola points).

(c) The proof uses only: at most two classes per line (Lagrange), the square arrangement of the four copies, and
the position of the partner classes $\sigma\kappa$, $\tau\kappa$ (Lemma~\ref{lem:partners}).  It does not use $p\ge7$;
$p=3$ and $p=5$ are covered as well.
\end{remark}

\section{Computer verification}

The script \texttt{slack/hjsw\_window\_check.py} of the repository recomputes, from the point set alone (no part of the
proof is trusted): all lines with at least three points of $\Pset$ (all have slope $\pm1$, three or four points, and lie
inside one orbit of classes); the identity $2|\Lc|+Z=3(p-1)$; the exact maximum by exhaustive search inside each orbit
(at most $2^{16}$ subsets), assembling a maximum witness which is passed through the independent verifier
\texttt{saturation.the\_law} in grid coordinates; and the number of maximum sets, compared with $9^{s}$.  All checks pass
for every prime $3\le p\le101$ (all $c\in\Fp^*$ for $p\le31$; $c\in\{1,2,3,-1,-2,(p+1)/2\}$ for larger $p$); log:
\texttt{bench/hjsw\_window\_check\_all.log}.  The values agree with the exact maxima obtained earlier by SAT
(kissat/CaDiCaL, two encodings, witnesses through the law): $18,30,36,48,54,66$ for $p=7,11,13,17,19,23$.

\begin{table}[h]\centering\small
\begin{tabular}{rrrrrrr}
\toprule
$p$ & $c$ & $|\Pset|$ & $|\Lc|$ & $Z$ & max & \#max sets\\
\midrule
7 & 1 & 24 & 8 & 2 & 18 & 9\\
11 & 1 & 40 & 14 & 2 & 30 & 9\\
13 & 1 & 48 & 16 & 4 & 36 & 81\\
13 & 2 & 48 & 18 & 0 & 36 & 1\\
17 & 1 & 64 & 22 & 4 & 48 & 81\\
19 & 1 & 72 & 26 & 2 & 54 & 9\\
23 & 1 & 88 & 32 & 2 & 66 & 9\\
101 & 1 & 400 & 148 & 4 & 300 & 81\\
\bottomrule
\end{tabular}
\caption{Sample values from the verification log ($|\Lc|=\tfrac32(p-1)-s$, $Z=2s$).}
\end{table}

\paragraph{Statement on authorship.}  The theorem, its proof and the verification program were produced by an
autonomous Claude agent (the third agent of the project) working from the setup and the computational observations
of the first two agents (exact maxima by SAT for $p\le23$); the owner set the question.  Everything is reproducible
from the repository.

\begin{thebibliography}{9}
\bibitem{HJSW75} R.~R. Hall, T.~H. Jackson, A.~Sudbery, K.~Wild, Some advances in the no-three-in-line problem,
 J. Combin. Theory Ser. A 18 (1975) 336--341.
\bibitem{KNS25} B.~Kov\'acs, Z.~L. Nagy, D.~R. Szab\'o, Randomised algebraic constructions for the no-$(k+1)$-in-line
 problem, arXiv:2508.07632 (2025); Advances in Combinatorics 2026:7.
\end{thebibliography}

\end{document}
